\documentclass[11pt]{article}
\usepackage[margin=1in]{geometry}
\usepackage{amsmath,amssymb,amsthm,mathtools}
\usepackage{booktabs,array}
\usepackage{graphicx}
\usepackage{natbib}
\usepackage{microtype}
\usepackage{setspace}
\usepackage{hyperref}
\usepackage{enumitem}
\usepackage{caption}
\hypersetup{colorlinks=true,citecolor=blue,linkcolor=blue,urlcolor=blue}
\captionsetup[figure]{name=Fig.,labelfont=bf,labelsep=space}
\setstretch{1.12}
\allowdisplaybreaks

\newtheorem{theorem}{Theorem}
\newtheorem{proposition}{Proposition}
\newtheorem{lemma}{Lemma}
\newtheorem{corollary}{Corollary}
\theoremstyle{remark}
\newtheorem{remark}{Remark}

\title{\textbf{Public Climate Protection and Private Production Redundancy}\\
\large Local Resilience Competition, Plant Location, and Disaster-State Market Power}
\author{}
\date{}

\begin{document}
\maketitle

\begin{abstract}
Climate-protection investments make production sites safer, but firms can respond by changing both primary locations and geographic backup capacity. This paper studies two jurisdictions that compete for mobile oligopolistic production through local protection against disruption. Firms choose primary production locations and costly backup readiness before localized disasters occur and surviving firms compete in a differentiated product market. Public protection directly lowers disruption risk but can crowd out private geographic redundancy. Their interaction creates a policy-ranking conflict: within the maintained symmetric primitive family, there is a nonempty open set in which coordinated welfare is maximized by zero additional local protection while decentralized jurisdictions choose a strictly positive symmetric protection level. A general marginal decomposition shows that the local incentive equals one quarter of the national marginal value plus a plant-attraction term. The result separates three margins relevant for climate-adaptation appraisal: the direct engineering benefit of protection, firms' endogenous private-resilience response, and interjurisdictional competition for production. Product-market rivalry matters because the profit difference between sole-survivor and duopoly states governs the private return to backup readiness. Thus place-based adaptation policy can alter the geographic organization of production as well as physical risk.
\end{abstract}

\noindent\textbf{Keywords:} climate adaptation; private resilience; firm location; public investment; oligopoly; fiscal competition.\\
\textbf{JEL:} H73, L13, Q54, R38.

\section{Related Literature}\label{sec:literature}

The paper connects three literatures. The contribution is the interaction among them, not a claim to introduce any one component.

First, a long literature studies strategic local public inputs and mobile firms. \citet{WalzWellisch1996} analyze strategic public provision of local inputs with endogenous oligopolistic location choice. \citet{MaurerWalz2000} study two regional governments competing for two mobile oligopolistic firms and show that regional competition generally produces suboptimal infrastructure provision. Related work shows that mobility and market integration can change the direction of local-public-good distortions \citep{Matsumoto1998,Ihara2008}. More recently, \citet{MengDingWang2024} document and model a negative relationship between local-government competition and industrial resilience operating through market segmentation and industrial-structure distortions. Our plant-attraction motive is deliberately close to this literature. The difference is that public protection also changes firms' privately chosen geographic redundancy, so the planner's marginal valuation of the same local input becomes endogenous to the industrial resilience response.

Second, environmental and regional models make firm location and market structure endogenous to policy. \citet{MarkusenMoreyOlewiler1993} show that environmental policy can induce discrete changes in plant location and market structure, and \citet{MarkusenMoreyOlewiler1995} study noncooperative regional environmental policy when plant locations are endogenous. Closest on strategic environmental policy and industrial allocation, \citet{MartinHerranEtAl2026} analyze two regions in a transboundary-pollution dynamic game in which environmental policy affects the long-run spatial distribution of firms. Their analysis emphasizes pollution policy, transport costs, and agglomeration/dispersion forces. Our mechanism instead treats public protection against disruption as the policy instrument and makes firms' costly geographic backup readiness endogenous; the resulting welfare channel runs through disaster-state production availability and market structure. \citet{OtazawaEtAl2016} incorporate natural-disaster risk into an economic-geography environment and show that market equilibrium may display excessive industrial agglomeration and vulnerability. Especially close to one part of our mechanism, \citet{GramesEtAl2019} study a firm's joint flood-risk location and protection decisions and show that existing flood defense can induce the firm to locate closer to water and invest less in private flood defense. We therefore do not claim novelty for policy-induced location change, adaptation crowd-out, or the joint occurrence of those two responses for a single firm.

Third, recent work studies private adaptation and resilience. \citet{Eisenack2014} analyzes private adaptation with endogenous market structure. \citet{CapponiDuStiglitz2024} show that market power and market incompleteness can produce insufficient supply-network resilience. \citet{CastroVincenziEtAl2024} study firms' geographic diversification of sourcing under climate risk. Most directly on public/private substitution, a 2026 working paper by Zhao, Yang, and Zhang (SSRN 6996377) reports evidence that public flood protection can crowd out private defensive investment and increase tail-risk losses. Accordingly, public/private resilience substitution is not our novelty claim.

The paper's contribution is narrower: a full-game policy-ranking result produced by combining protection-induced private geographic-redundancy adjustment, mobile oligopolistic plant location, and decentralized competition for plants. In the fixed-redundancy benchmark, the model resembles local-public-input competition. In the fixed-location benchmark, public/private resilience substitution remains but the plant-attraction motive disappears. In the full game, Proposition~\ref{prop:attraction-threshold} shows transparently how the local plant-attraction term can overturn a negative national marginal value of additional protection, while Theorem~\ref{thm:main} establishes a global policy-ranking reversal on a nonempty open set. This policy-ranking reversal, rather than any single ingredient, is the paper's novelty claim.

\section{Model}\label{sec:model}

There are two jurisdictions, $A$ and $B$, two ex ante symmetric firms, and a representative consumer. Jurisdictions choose place-based climate protection, firms choose primary production locations and geographic backup readiness, localized disasters are realized, and surviving firms compete in the product market. The model is deliberately static. Its purpose is to isolate one interaction: local competition for mobile production changes public protection, while public protection changes firms' private production redundancy.

\subsection{Product market}

Consumer utility is
\begin{equation}
U(x_1,x_2)=x_1+x_2-\frac12\left(x_1^2+x_2^2+2\gamma x_1x_2\right),
\qquad 0\leq \gamma<1,
\label{eq:utility}
\end{equation}
where $x_i$ is consumption of firm $i$'s product and $\gamma$ indexes product substitutability. The corresponding inverse demands are
\begin{equation}
p_i=1-x_i-\gamma x_j.
\label{eq:inverse-demand}
\end{equation}
Marginal production cost is normalized to zero. Whenever a firm is available through either its primary plant or its backup arrangement, it can supply the unconstrained Cournot quantity at this same marginal cost; there is no separate state-dependent capacity constraint. If both firms are available after the disaster state is realized, they compete in quantities. If only one firm is available, it is a monopolist. If neither firm is available, output is zero.

\subsection{Public protection and localized risk}

Jurisdiction $j\in\{A,B\}$ chooses public climate protection
\begin{equation}
a_j\in[0,\bar a].
\end{equation}
Protection lowers the probability that primary production located in $j$ is disrupted:
\begin{equation}
q_j=q_0-a_j,
\qquad 0<\bar a<q_0<1.
\label{eq:risk}
\end{equation}
The resource cost of protection is $c a_j^2/2$, with $c>0$. The linear risk technology defines the baseline theorem environment. Section~\ref{sec:discussion} reports a diagnostic portability check under a decreasing nonlinear risk technology; that check is robustness evidence rather than an extension of the exact theorem to every decreasing risk function.

Localized shocks generate correlation in primary production failures. Let $F_i\in\{0,1\}$ indicate failure of firm $i$'s primary production. If both firms place their primary plants in the same jurisdiction $j$, their primary plants fail jointly with probability $J=q_j$. If they locate in different jurisdictions, regional shocks are independent and joint primary failure occurs with probability
\begin{equation}
J=q_Aq_B.
\label{eq:joint-risk}
\end{equation}
Thus $\Pr(F_i=1)=s_i$ and $\Pr(F_1=F_2=1)=J$, with the location configuration determining $s_i$ and $J$.

\subsection{Primary location and geographic backup readiness}

Firm $i$ chooses a primary production location $L_i\in\{A,B\}$. Let
\begin{equation}
s_i=q_{L_i}
\end{equation}
be the failure probability of its primary production. After primary locations are chosen, firm $i$ chooses backup readiness $r_i\in[0,1]$. Conditional on primary failure, backup production is available with probability $r_i$. Maintaining readiness costs
\begin{equation}
C_i^R(r_i)=\frac{k}{2}r_i^2,
\qquad k>0,
\label{eq:backup-cost}
\end{equation}
and this cost is paid ex ante and is state independent.

The joint law of backup success is a maintained primitive. Conditional on public policies, primary locations, and the realized primary-failure vector $(F_1,F_2)$, backup-success draws are independent across firms. Whenever $F_i=1$, firm $i$'s backup is available with conditional probability $r_i$; conditional backup success does not otherwise depend on the regional shock that generated the primary-failure vector. The backup-success draws are also independent of the location-fit shocks introduced below. Equivalently, firm $i$ is unavailable with probability $s_i(1-r_i)$, while both firms are unavailable with probability
\begin{equation}
J(1-r_1)(1-r_2).
\label{eq:joint-unavailability}
\end{equation}
This conditional-independence assumption is substantive: fixing only the marginal success probabilities $r_i$ would not determine the joint production-availability states. The reduced-form variable $r_i$ should therefore be read as the success probability of the firm's geographically separate contingency arrangement after its primary site fails; any hazard exposure specific to that arrangement is already incorporated into $r_i$, and there is no additional common backup shock in the baseline model.

This reduced-form readiness choice captures costly alternative facilities, pre-arranged production capacity, or other geographically separate production options. It is smooth by construction; none of the main results depends on crossing a fixed-cost threshold.

To avoid a knife-edge winner-take-all location response, each firm receives a privately observed idiosyncratic location-fit shifter favoring $A$ relative to $B$,
\begin{equation}
\varepsilon_i\sim U[-H,H],\qquad H>0,
\label{eq:location-shock}
\end{equation}
independently across firms and independently of all regional-disaster and backup-success draws. After public policies are chosen and before the simultaneous location decisions, firm $i$ privately observes $\varepsilon_i$. Each firm knows the distribution of the rival's shifter and forms beliefs over the rival's location choice in the usual Bayesian way. We use $\varepsilon_i$ as a smoothing device for the discrete location decision, not as consumption, accounting profit, or a real resource. It therefore affects the firm's cutoff rule but is not part of baseline material social surplus. Given continuation profits, firm $i$ chooses $A$ when its deterministic payoff advantage of $A$ plus $\varepsilon_i$ is nonnegative. Appendix~\ref{app:verification} also shows that counting the selected fit payoff in welfare leaves the canonical policy ranking unchanged.

\subsection{Governments and welfare}

Let $\mathcal S(a_A,a_B)$ denote expected national product-market surplus net of private backup costs, after location and backup choices are resolved. A primary plant located in a jurisdiction generates local employment, tax-base, or immobile-factor benefits summarized by $b>0$. The baseline incidence rule assigns one half of national market surplus $\mathcal S$ to each jurisdiction. If $p$ denotes the symmetric probability that a firm chooses $A$, jurisdiction $A$'s payoff is
\begin{equation}
G_A=\frac12\mathcal S+2bp-\frac{c}{2}a_A^2,
\label{eq:government-payoff}
\end{equation}
and $B$ is symmetric. This equal-share rule is a maintained reduced-form incidence assumption, not an accounting identity derived from the product market. A coordinated planner values the same real product-market surplus and both local plant benefits:
\begin{equation}
W=\mathcal S+2b-\frac{c}{2}(a_A^2+a_B^2).
\label{eq:national-welfare}
\end{equation}
Because there are always two primary plants, $2b$ is constant at the national level. With the equal-share incidence rule, $G_A+G_B=W$. The decentralized governments therefore have a plant-attraction motive that is absent from the planner's location accounting.

\subsection{Timing and equilibrium concept}

The game proceeds as follows.
\begin{enumerate}[label=\arabic*.]
\item Jurisdictions simultaneously choose $(a_A,a_B)$.
\item Each firm privately observes its own $\varepsilon_i$; firms then simultaneously choose primary locations, forming beliefs over the rival's private draw.
\item Firms simultaneously choose backup readiness $(r_1,r_2)$ after primary locations are observed.
\item Localized primary-disaster states are realized; conditional on primary failures, backup success is realized according to the conditional-independence law above.
\item Available firms compete in quantities.
\end{enumerate}
We solve by backward induction. The location stage is a symmetric Bayesian Nash equilibrium in cutoff strategies. The analysis uses the interior continuation region certified for the stated policy domain. Upstream deviations are evaluated by re-solving the downstream location and backup subgames rather than holding on-path continuation formulas fixed.

\section{Equilibrium Characterization}\label{sec:equilibrium}

\subsection{Product-market payoffs}

When both firms survive, the Cournot equilibrium is
\begin{equation}
x_D=p_D=\frac{1}{2+\gamma},
\qquad
\pi_D=\frac{1}{(2+\gamma)^2}.
\label{eq:duopoly}
\end{equation}
When only one firm survives,
\begin{equation}
x_M=p_M=\frac12,
\qquad
\pi_M=\frac14.
\label{eq:monopoly}
\end{equation}
Define the disaster-state scarcity-rent wedge
\begin{equation}
\Delta(\gamma)\equiv \pi_M-\pi_D
=\frac{\gamma(\gamma+4)}{4(\gamma+2)^2}.
\label{eq:delta}
\end{equation}
For $\gamma>0$, $\Delta>0$, and
\begin{equation}
\Delta'(\gamma)=\frac{2}{(\gamma+2)^3}>0.
\label{eq:delta-derivative}
\end{equation}
Thus closer substitutes increase the profit gain from being the sole surviving producer relative to competing in the disaster state.

\subsection{Backup choice}

Fix primary locations and public policies. Let $s_i$ and $s_j$ be the firms' primary failure probabilities and let $J$ be their joint primary failure probability. After integrating over disaster states, firm $i$'s expected profit can be written as
\begin{equation}
\Pi_i
=\pi_D-\pi_D s_i(1-r_i)
+\Delta s_j(1-r_j)
-\Delta J(1-r_i)(1-r_j)
-\frac{k}{2}r_i^2.
\label{eq:expected-profit}
\end{equation}
The interior first-order condition is
\begin{equation}
kr_i=\pi_D s_i+\Delta J(1-r_j).
\label{eq:backup-foc}
\end{equation}
The own second derivative is $-k<0$, so the best response is unique conditional on the rival's readiness.

\begin{proposition}[Strategic substitution in backup readiness]\label{prop:strategic-sub}
On the interior continuation branch,
\begin{equation}
\frac{\partial r_i^{BR}}{\partial r_j}=-\frac{\Delta J}{k}<0
\end{equation}
whenever $\gamma>0$ and $J>0$.
\end{proposition}

The economic force is specific to competition in the disruption state. A rival that remains productive after a common disruption eliminates the monopoly component of the firm's contingency payoff and therefore lowers its incentive to maintain backup readiness. Proposition~\ref{prop:strategic-sub} is a mechanism lemma, not a novelty claim.

For symmetric public risk $q$, the two location configurations deliver especially transparent expressions. If primary plants are co-located,
\begin{equation}
r^C(q)=\frac{q\pi_M}{k+q\Delta}.
\label{eq:rC}
\end{equation}
If primary plants are geographically dispersed,
\begin{equation}
r^D(q)=\frac{q(\pi_D+q\Delta)}{k+q^2\Delta}.
\label{eq:rD}
\end{equation}
Both are increasing in $q$, hence decreasing in public protection. Both are also decreasing in $\gamma$:
\begin{equation}
\frac{\partial r^C}{\partial \gamma}<0,
\qquad
\frac{\partial r^D}{\partial \gamma}<0.
\label{eq:gamma-r}
\end{equation}
The latter comparative static follows because a larger scarcity-rent wedge makes the value of backup more dependent on the rival's failure rather than on preserving competitive production in every state.

\subsection{Primary location}

Let $\Pi^{XY}_i$ denote firm $i$'s equilibrium continuation profit when firm $1$ locates in $X\in\{A,B\}$ and firm $2$ in $Y\in\{A,B\}$. For firm 1, define
\begin{align}
d_A&\equiv \Pi^{AA}_1-\Pi^{BA}_1,\\
d_B&\equiv \Pi^{AB}_1-\Pi^{BB}_1.
\end{align}
If the rival chooses $A$ with probability $p$, the deterministic payoff difference between choosing $A$ and $B$ is
\begin{equation}
D(p)=p d_A+(1-p)d_B.
\end{equation}
Because $\varepsilon_i$ is privately observed, the firm's cutoff rule implies the affine Bayesian best response
\begin{equation}
BR(p)=\frac12+\frac{D(p)}{2H}
=u_0+Bp,
\label{eq:location-br}
\end{equation}
where
\begin{equation}
u_0\equiv\frac{H+d_B}{2H},
\qquad
u_1\equiv\frac{H+d_A}{2H},
\qquad
B\equiv\frac{d_A-d_B}{2H}.
\label{eq:location-endpoints}
\end{equation}
Here $u_0=BR(0)$ and $u_1=BR(1)$. Whenever $u_0,u_1\in(0,1)$, the affine expression is valid globally on $p\in[0,1]$: no probability clipping occurs at either endpoint or anywhere between them. If in addition $|B|<1$, $BR$ is a contraction and the simultaneous private-information location game has a unique interior probability fixed point. Solving $p=BR(p)$ gives
\begin{equation}
p=\frac{H+d_B}{2H-d_A+d_B}.
\label{eq:location-prob}
\end{equation}
At the canonical witness, the exact whole-domain certificate verifies $0<u_0,u_1<1$ and $|B|<1$ for every $(a_A,a_B)\in[0,\bar a]^2$. Thus equation~\eqref{eq:location-prob} is the unique location continuation after every admissible first-stage deviation, not merely a local interior solution. At symmetric public policies $a_A=a_B$, symmetry implies $p=1/2$ and $D=0$. Off the symmetric policy path, equation~\eqref{eq:location-prob} captures the plant-attraction effect of a unilateral protection increase.

\section{Local Protection Competition and the Main Result}\label{sec:main}

The preceding section isolates two responses to public protection. First, a safer jurisdiction becomes more attractive as a primary production location. Second, lower primary risk reduces firms' willingness to maintain geographically separate backup capacity. This section shows that the interaction of those responses can reverse the social policy ranking.

\subsection{The two marginal incentives}

At a symmetric policy profile, define the planner's marginal payoff from a common increase in public protection by
\begin{equation}
M_S(\gamma)\equiv \left.\frac{dW(a,a)}{da}\right|_{a=0},
\label{eq:MS}
\end{equation}
and define jurisdiction $A$'s marginal payoff from a unilateral increase by
\begin{equation}
M_L(\gamma)\equiv \left.\frac{\partial G_A(a_A,0)}{\partial a_A}\right|_{a_A=0}.
\label{eq:ML}
\end{equation}
The planner internalizes the change in product-market surplus and private redundancy. A local government internalizes only its assigned share of national market surplus but additionally values the increase in its expected primary-plant count.

For the rational parameter vector
\begin{equation}
\theta_0=
\left(
\gamma,q_0,k,H,b,c,\bar a
\right)
=
\left(
\frac{12}{25},\frac3{10},\frac{169}{3000},\frac1{100},\frac{11}{200},3,\frac2{25}
\right),
\label{eq:witness}
\end{equation}
exact rational evaluation gives
\begin{equation}
M_S(12/25)<0<M_L(12/25).
\label{eq:marginal-conflict}
\end{equation}
Thus, at the same policy profile, coordinated welfare falls when both jurisdictions add protection, while a jurisdiction gains from adding protection unilaterally.

\begin{theorem}[Protection--Attraction Conflict]\label{thm:main}
There exists a nonempty open set of primitives for which coordinated welfare is maximized by zero additional local climate protection while the decentralized jurisdictional game has a strictly positive symmetric protection equilibrium. At the canonical rational witness in \eqref{eq:witness}, the coordinated optimum is $(a_A,a_B)=(0,0)$ and the positive symmetric decentralized equilibrium satisfies
\begin{equation}
a^{LG}\in(0.02301968485,\;0.02301968495).
\label{eq:alg-interval}
\end{equation}
All material downstream continuation equilibria used to evaluate first-stage deviations remain interior on the frozen policy domain.
\end{theorem}

The proof is computer-assisted but exact at the global-policy step. Exact rational root isolation certifies the unique symmetric local-government first-order-condition root. Tensor-product Bernstein certificates over the full policy rectangle prove that coordinated welfare is strictly decreasing in each protection coordinate, that all continuation objects remain regular and interior, and that the local government's objective is strictly concave in its own policy when the rival chooses the isolated equilibrium root. Hence the isolated stationary point is a global best response and the positive symmetric profile is a Nash equilibrium. Numerical global-deviation routines are retained only as independent stress tests. Appendix~\ref{app:verification} documents the exact and numerical layers separately.

\subsection{A product-substitutability interval}

Holding $(q_0,k,H,b,c,\bar a)$ at the witness values and varying $\gamma$, exact root isolation identifies a local-government threshold
\begin{equation}
\gamma_L\in(0.4745893334177,\;0.4745893334186)
\end{equation}
and a planner threshold
\begin{equation}
\gamma_S\in(0.4982912217039,\;0.4982912217048).
\end{equation}
Hence $\gamma_L<\gamma_S$. For $\gamma$ between these thresholds, the local marginal incentive at zero is positive while the coordinated marginal incentive is negative. Figure~\ref{fig:policy-regime} shows the corresponding symmetric policy comparison. The figure is an exposition device generated from the frozen model; the theorem itself rests on the exact witness and verification record.

\begin{figure}[t]
\centering
\includegraphics[width=.78\textwidth]{../figures/policy_regime.pdf}
\caption{Public protection and product substitutability. The vertical lines mark the exact root-isolation thresholds $\gamma_L$ and $\gamma_S$. The plotted policies are generated from the frozen symmetric-policy equations; they are not additional calibrated moments.}
\label{fig:policy-regime}
\end{figure}

\subsection{Economic interpretation}

The theorem is not a general statement that climate protection is inefficient. The conflict arises in a specific margin of place-based industrial resilience. Public protection has an engineering benefit: it lowers the probability that a primary production site fails. But it also changes firms' organizational response. Safer primary production reduces the private value of maintaining costly geographic backup capacity. Once that endogenous reduction in backup capacity makes an additional common protection increment socially undesirable, a coordinated planner stops. A local government need not stop because the same protection increment attracts primary production from the competing jurisdiction. The decentralized policy therefore remains positive even though its national marginal value has turned negative.

\section{Welfare and the Public--Private Resilience Substitution}\label{sec:welfare}

\subsection{Exact surplus accounting}

From utility \eqref{eq:utility}, consumer surplus when both firms survive is
\begin{equation}
CS_D=\frac{1+\gamma}{(2+\gamma)^2}.
\end{equation}
With total producer surplus $2/(2+\gamma)^2$, total surplus in the duopoly state is
\begin{equation}
T_D=\frac{3+\gamma}{(2+\gamma)^2}.
\label{eq:TD}
\end{equation}
When only one firm survives,
\begin{equation}
CS_M=\frac18,\qquad PS_M=\frac14,\qquad T_M=\frac38.
\label{eq:TM}
\end{equation}

Let
\begin{equation}
m_i=s_i(1-r_i)
\end{equation}
be firm $i$'s probability of being unavailable. Under the conditional-independence backup-success law in Section~\ref{sec:model}, the probability that both firms are unavailable is
\begin{equation}
z=J(1-r_1)(1-r_2).
\end{equation}
The four final availability probabilities are therefore $1-m_1-m_2+z$, $m_2-z$, $m_1-z$, and $z$. Expected product-market surplus net of backup costs is
\begin{equation}
\mathcal S
=T_D+(T_M-T_D)(m_1+m_2)+(T_D-2T_M)z
-\frac{k}{2}(r_1^2+r_2^2).
\label{eq:surplus}
\end{equation}
The welfare comparison therefore reflects real output losses, consumer surplus, producer surplus, and real public and private resilience costs. It is not generated by treating fiscal transfers as social costs.

\subsection{Channel decomposition at the canonical witness}

The central welfare logic can be seen at the no-additional-protection point. Consider a symmetric common increase in $a$ at the witness $\gamma=12/25$. Holding firms' backup readiness fixed, the direct reduction in production risk raises coordinated welfare. Allowing backup readiness to respond reverses the sign. The generated decomposition is reported in Table~\ref{tab:channels}.

\begin{table}[t]
\centering
\caption{Marginal welfare channels at the canonical witness}
\label{tab:channels}
\begin{tabular}{lr}\toprule
Component & Marginal effect \\\midrule
Direct protection effect, redundancy fixed & 0.04146 \\
Endogenous redundancy response & -0.04794 \\
Total coordinated effect & -0.00648 \\
Local-government unilateral payoff effect & 0.00193 \\\bottomrule
\end{tabular}

\end{table}

The direct protection effect is approximately $+0.04146$. The endogenous redundancy response contributes approximately $-0.04794$, leaving a coordinated marginal effect near $-0.00648$. At the same point, a jurisdiction's unilateral marginal payoff is approximately $+0.00193$. The signs, rather than the magnitudes, are the economically relevant objects. The direct protection benefit is positive; the induced reduction in private geographic redundancy is sufficiently strong to make the social total negative; and plant attraction keeps the decentralized marginal incentive positive.

\subsection{Why product-market competition matters}

The scarcity-rent wedge $\Delta(\gamma)$ links the industrial-organization block to the resilience decision. A higher $\gamma$ lowers duopoly profit relative to monopoly profit, thereby increasing $\Delta$. Equations \eqref{eq:rC}--\eqref{eq:rD} then imply that equilibrium backup readiness falls with product substitutability on the maintained symmetric interior branch. This does not mean that competition always lowers resilience in richer environments, nor does it assign a general sign to the interaction between product substitutability and the marginal response of readiness to protection. In the maintained baseline model, backup readiness is valuable partly because it determines who remains active after a localized disruption, so the relative value of competitive survival versus sole survival is an essential determinant of the firm's private resilience incentive.

\section{Nested Benchmarks}\label{sec:benchmarks}

The contribution depends on an interaction, so the appropriate comparison is not whether each component is individually familiar. It is whether the mechanism in Theorem~\ref{thm:main} remains available after either strategic margin is removed. The benchmarks below are nested interventions on the full game, not additional global-equilibrium theorems.

In the \emph{fixed-primary-location} benchmark, the primary-location profile is assigned exogenously before the policy comparison and does not respond to any policy deviation; the backup game, disaster-state probabilities, and product-market continuation are otherwise solved exactly as in the baseline model. In the \emph{fixed-private-redundancy} benchmark, readiness levels are treated as exogenous constants during the policy comparison and do not respond to protection; the location stage, primary-risk mapping, and product-market continuation remain active. Thus each benchmark removes exactly one endogenous response while retaining the other baseline primitives. Table~\ref{tab:benchmarks} summarizes the logic.

\begin{table}[t]
\centering
\caption{Nested benchmark logic}
\label{tab:benchmarks}
\small
\begin{tabular}{p{.20\textwidth}p{.25\textwidth}p{.46\textwidth}}
\toprule
Model & Endogenous margins & Policy consequence \\
\midrule
Fixed primary locations & Private redundancy only & Public/private substitution remains, but protection cannot attract primary production. \\
Fixed private redundancy & Primary location only & Standard local public-input competition remains, but public policy cannot reorganize private resilience. \\
Full model & Primary location and private redundancy & Plant attraction operates after redundancy crowd-out changes the social value of protection; $a^{LG}>0=a^{SP}$ can arise. \\
\bottomrule
\end{tabular}
\end{table}

\subsection{Fixed primary locations}

Fix any primary-location profile and suppress the location stage. A unilateral protection increase then cannot attract a primary plant: the expected local plant count is constant and the derivative of the $2bp$ term in \eqref{eq:government-payoff} with respect to local protection is zero. The backup game is still re-solved after every policy change, so public protection can continue to crowd out private readiness. This benchmark therefore isolates public/private resilience substitution without the plant-attraction channel.

\subsection{Fixed private redundancy}

Fix the firms' readiness levels exogenously and hold them unchanged under every policy deviation. Primary locations are still chosen through the Bayesian location game, and protection still changes primary disruption risk and location attractiveness. The model therefore retains competition for mobile production through a local public input, but public policy can no longer reorganize firms' contingency capacity. This benchmark isolates the plant-attraction channel without endogenous resilience adjustment.

\subsection{Full-game interaction}

The full model combines the two margins but does more than place them side by side. Public protection changes the private resilience architecture, which changes the national value of additional protection. Local governments then act on a different objective because they value the induced relocation of primary production. The canonical calculations verify the direction of each benchmark mechanism, while Theorem~\ref{thm:main} and Appendix~\ref{app:verification} certify the global policy-ranking reversal only for the full game. The benchmark comparison is therefore a mechanism-identification exercise, not a claim that each restricted game separately satisfies the full theorem's global policy ranking.

\section{Institutional Interpretation and Empirical Implications}\label{sec:institutions}

The model is written as a theory of manufacturing production under localized climate risk, but its primitives correspond to recognizable policy and business-continuity decisions. The institutional evidence is used to discipline interpretation, not to claim that the exact strategic channel has already been identified empirically.

\subsection{Local climate adaptation and business continuity}

Japan's Ministry of the Environment explicitly assigns local governments a role in regional climate-adaptation planning and, in its 2026 revision of the local adaptation-plan manual, discusses progress management and joint plan formation by multiple local governments \citep{MOE2026}. This supports the model's treatment of place-based public adaptation as a decentralized policy choice.

Firms also make real geographic continuity choices. The Small and Medium Enterprise Agency's business-continuity guidance lists alternative facilities, parallel operation at another site, temporary replacement facilities, inter-firm arrangements, and construction of facilities at another location as possible responses to disruption \citep{SMEA_BCP}. These are direct institutional counterparts to the model's reduced-form backup-readiness choice. The model does not require every firm literally to own a duplicate plant; it requires costly preparations that preserve alternative production capability when a primary site fails.

\subsection{Competition for mobile production}

The Ministry of Economy, Trade and Industry maintains industrial-site information, supports industrial-site development, and facilitates matching between firms and candidate locations \citep{METI2026}. Local governments also explicitly pursue business attraction for employment and regional economic activity. For example, Ehime Prefecture describes enterprise attraction as a policy for employment expansion and regional economic revitalization \citep{EHIME2026}. These facts support the $b$ term in \eqref{eq:government-payoff}: primary production has locally valued benefits that need not be additional national benefits when activity is relocated across jurisdictions.

The evidence does \emph{not} establish that local climate-protection spending is currently chosen for the purpose of attracting firms. That direct causal link is an implication to be tested, not an assumed fact.

\subsection{Empirical predictions}

The theory yields several testable implications, with the scope of each comparative static stated explicitly.
\begin{enumerate}[label=(\roman*)]
\item Improvements in place-based protection should increase primary production exposure to the protected jurisdiction while reducing firms' use of geographically separate contingency capacity, other things equal.
\item On the maintained symmetric interior branch, industries with closer substitutes have lower equilibrium backup readiness because the sole-survivor versus duopoly profit wedge is larger. This is a \emph{level} prediction. The model does not generally predict that product substitutability strengthens the marginal crowd-out effect of public protection. At the canonical co-located symmetric origin, the exact cross-partial is negative (Appendix~\ref{app:verification}), so the readiness reduction induced by an incremental increase in protection is locally \emph{weaker}, not stronger, when products are closer substitutes at this point.
\item Jurisdictions placing a larger value on plant attraction---for example because local employment or immobile-factor rents are more salient---should exhibit a larger decentralized protection incentive relative to a coordinated benchmark.
\item Evaluation of regional resilience projects based only on engineering reductions in local failure probability may overstate net industrial resilience benefits if firms respond by reducing costly alternative production arrangements.
\end{enumerate}
These predictions distinguish the mechanism from a pure safe-development story in which only the exposure of households or physical capital changes after public protection. Prediction (ii) deliberately separates the level effect of product substitutability from the interaction between substitutability and the marginal policy response; the latter is not assigned a global sign by the model.

\section{Related Literature}\label{sec:literature}

The paper connects three literatures. The contribution is the interaction among them, not a claim to introduce any one component.

First, a long literature studies strategic local public inputs and mobile firms. \citet{WalzWellisch1996} analyze strategic public provision of local inputs with endogenous oligopolistic location choice. \citet{MaurerWalz2000} study two regional governments competing for two mobile oligopolistic firms and show that regional competition generally produces suboptimal infrastructure provision. Related work shows that mobility and market integration can change the direction of local-public-good distortions \citep{Matsumoto1998,Ihara2008}. Our plant-attraction motive is deliberately close to this literature. The difference is that public protection also changes firms' privately chosen geographic redundancy, so the planner's marginal valuation of the same local input becomes endogenous to the industrial resilience response.

Second, environmental and regional models make firm location and market structure endogenous to policy. \citet{MarkusenMoreyOlewiler1993} show that environmental policy can induce discrete changes in plant location and market structure, and \citet{MarkusenMoreyOlewiler1995} study noncooperative regional environmental policy when plant locations are endogenous. \citet{OtazawaEtAl2016} incorporate natural-disaster risk into an economic-geography environment and show that market equilibrium may display excessive industrial agglomeration and vulnerability. We therefore do not claim novelty for policy-induced location change or for the observation that agglomeration may create disaster exposure. Our organizational margin is instead firm-level backup readiness chosen after primary location.

Third, recent work studies private adaptation and resilience. \citet{Eisenack2014} analyzes private adaptation with endogenous market structure. \citet{CapponiDuStiglitz2024} show that market power and market incompleteness can produce insufficient supply-network resilience. \citet{CastroVincenziEtAl2024} study firms' geographic diversification of sourcing under climate risk. Most directly, \citet{ZhaoYangZhang2026} provide evidence that public flood protection can crowd out private defensive investment and increase tail-risk losses. Accordingly, public/private resilience substitution is not our novelty claim.

The paper's result is a full-game interaction. In the fixed-redundancy benchmark, the model resembles local-public-input competition. In the fixed-location benchmark, public/private resilience substitution remains but the plant-attraction motive disappears. Only the full game can generate a regime in which private redundancy adjustment makes additional common public protection socially undesirable while a local government nevertheless supplies positive protection to attract mobile oligopolistic production.

\section{Discussion}\label{sec:discussion}

\subsection{What the result does and does not say}

The model does not imply that real-world flood protection, sea walls, drainage systems, or other adaptation investments are generally excessive. Such projects protect households, lives, public infrastructure, nonindustrial assets, and sectors outside the model. The analysis isolates an industrial-organization margin that standard project appraisal can miss: firms may reduce private geographic backup capacity in response to stronger place-based public protection, and decentralized governments may still value the same protection as a tool for attracting production.

The appropriate policy interpretation is therefore joint evaluation. A regional protection project that materially changes firms' location and continuity decisions should be evaluated together with those induced private organizational responses. The theorem identifies circumstances in which the marginal industrial-resilience component of a project can become negative even though the engineering reliability effect remains positive.

\subsection{Why the model is kept minimal}

Several extensions are natural but intentionally excluded from the baseline. Firms could choose multiple backup sites, governments could offer tax subsidies, disasters could be spatially correlated across jurisdictions, or the product market could use price rather than quantity competition. These may matter for quantitative applications, but none is required to state the present mechanism. Adding them before establishing the benchmark would obscure the distinction between the paper's contribution and established results on fiscal competition, adaptation crowd-out, or resilient supply networks.

The reduced-form backup choice should likewise not be read as a literal probability purchased from a technology. It is a tractable index of costly preparedness that increases the probability of maintaining production after primary-site disruption. The institutional examples in Section~\ref{sec:institutions} show that firms have multiple practical margins that fit this interpretation.

\subsection{Generality}

The same strategic logic can apply outside manufacturing whenever four ingredients are present: location-specific correlated disruption risk, public investment that lowers local risk, mobile competing firms, and costly private geographic redundancy. Data centers with geographic failover capacity and logistics operators with alternative terminal or routing capacity are possible examples. These interpretations require no new strategic force, but the paper keeps manufacturing as its canonical application.

\section{Conclusion}\label{sec:conclusion}

Public climate protection and private production resilience are jointly determined. A safer jurisdiction can attract primary production, but the same protection can reduce firms' willingness to maintain costly geographic backup capacity. When local governments compete for mobile production, they value the relocation benefit even if the induced reduction in private redundancy has already made additional public protection socially undesirable.

A simple two-jurisdiction oligopoly model formalizes this interaction. Firms choose primary locations and backup readiness before localized disasters, and surviving firms compete in the product market. The model generates a nonempty open set of primitives in which the coordinated planner chooses no additional local protection while decentralized jurisdictions choose strictly positive protection. The policy conflict is absent from the relevant nested benchmarks: fixed locations remove the attraction motive, while fixed redundancy removes the organizational response that changes the planner's valuation of protection.

The broader implication is methodological. Place-based resilience policy should not be evaluated as an engineering change in local hazard exposure holding private organization fixed. Where firms can move production and redesign contingency capacity, public adaptation changes the industrial resilience architecture itself. Decentralization then adds a second wedge because local governments value where production is located, not only whether the national production system remains resilient.


\section*{Data and Code Availability}
This is a theoretical study and uses no empirical dataset. The computer-assisted proof code, exact certificates, generated outputs, and instructions needed to reproduce the reported results are supplied to reviewers in an anonymized replication archive. Following successful peer review and before final acceptance, the replication materials will be deposited in a public repository and the persistent repository information will be supplied to the journal.

\section*{AI Assistance Disclosure}
OpenAI generative-AI tools were used during manuscript development for research assistance, adversarial mathematical and code review, and language editing. All mathematical claims, computer-assisted certificates, references, and generated outputs were independently checked against the stated model primitives and reproducibility tests. The author retains full responsibility for the scholarly content and final manuscript. No AI system is listed as an author.

\appendix
\section{Proofs and Verification}\label{app:verification}

This appendix records the analytical identities used in the paper and the computer-assisted verification protocol for the main existence result. The reproducibility package separates symbolic certificates, numerical continuation audits, and an independent state-enumeration payoff evaluator.

\subsection{Product-market and backup results}

The Cournot first-order conditions from \eqref{eq:inverse-demand} are
\begin{equation}
1-2x_i-\gamma x_j=0.
\end{equation}
Symmetry yields $x_i=1/(2+\gamma)$ and \eqref{eq:duopoly}; the monopoly result \eqref{eq:monopoly} is immediate. Subtracting profits gives \eqref{eq:delta}, and direct differentiation gives \eqref{eq:delta-derivative}.

To obtain \eqref{eq:expected-profit}, enumerate four production-availability states after primary and backup outcomes. Differentiating with respect to $r_i$ gives
\begin{equation}
\frac{\partial\Pi_i}{\partial r_i}
=\pi_D s_i+\Delta J(1-r_j)-kr_i,
\end{equation}
which yields \eqref{eq:backup-foc}. The cross-best-response derivative in Proposition~\ref{prop:strategic-sub} follows directly. Under symmetry, substituting $J=q$ for co-location and $J=q^2$ for dispersion gives \eqref{eq:rC} and \eqref{eq:rD}. Their derivatives with respect to $q$ are positive on the frozen interior region. Differentiation with respect to $\Delta$, combined with $\Delta'(\gamma)>0$, gives \eqref{eq:gamma-r}.

\subsection{Exact threshold certificates}

At the rational values $(q_0,k,H,b,c,\bar a)=(3/10,169/3000,1/100,11/200,3,2/25)$, the marginal objects $M_S(\gamma)$ and $M_L(\gamma)$ are rational functions of $\gamma$. Their exact numerator and denominator coefficients are stored in the reproducibility artifact \texttt{docs/certificate\_polynomials.json}. Exact real-root isolation establishes that the relevant numerator root of $M_L$ lies in
\begin{equation}
(0.4745893334177016,\;0.4745893334185960),
\end{equation}
and that the relevant numerator root of $M_S$ lies in
\begin{equation}
(0.4982912217039472,\;0.4982912217047885).
\end{equation}
The denominators have no root on the maintained $\gamma$ interval. Hence the ordering $\gamma_L<\gamma_S$ is exact, not a finite-difference pattern.

At $\gamma=12/25$, the numerator polynomial of the symmetric local-government first-order condition has a unique root on $[0,2/25]$, isolated in the interval reported in \eqref{eq:alg-interval}. The symmetric planner derivative is negative throughout the corresponding symmetric policy domain; its exact derivative numerator has no root on that interval and is negative at zero.

\subsection{Whole-domain continuation audit}\label{app:continuation}

The first-stage game is sequential. A government deviation changes local failure risk, which changes location incentives and the backup subgame. The verification therefore re-solves the downstream location fixed point and backup Nash equilibrium after each material government deviation rather than substituting the deviation into on-path formulas.

At the canonical witness, the independent numerical audit searches $a_A\in[0,\bar a]$ against the rival's candidate equilibrium policy using 2,001 policy points and bounded optimization. Throughout this domain, location probabilities and backup readiness remain strictly interior. The audit records zero unresolved continuation calls and zero numerical failures. A separate two-dimensional planner audit evaluates the full $(a_A,a_B)\in[0,\bar a]^2$ domain on a deterministic grid and from multiple continuous-optimization starting points; all checks select $(0,0)$ at the witness.

The independent evaluator enumerates the four production-availability states directly from primitive probabilities and reconstructs expected profit and surplus without calling the candidate closed-form payoff routine. It agrees with the closed form at the maintained test points. This separation is intended to prevent a common-code error from serving simultaneously as derivation and verification.

\subsection{Open-set argument}

At the witness, the planner and local-government marginal inequalities are strict; the positive symmetric government root is simple and interior; location probabilities and backup choices are strictly interior; and the continuation mapping remains within the verified regular region. These objects are continuous in primitives. Hence the sign conflict and positive local equilibrium persist in a sufficiently small neighborhood of the witness. This gives the open-set part of Theorem~\ref{thm:main}. The theorem is therefore an existence result, not a claim that every parameter configuration generates excessive local protection.


\begin{thebibliography}{99}
\bibitem[Capponi et~al.(2024)Capponi, Du, and Stiglitz]{CapponiDuStiglitz2024}
Capponi, A., Du, C., and Stiglitz, J.~E. (2024).
\newblock Are supply networks efficiently resilient?
\newblock \emph{Finance and Economics Discussion Series} 2024-031, Board of Governors of the Federal Reserve System. doi:10.17016/FEDS.2024.031.

\bibitem[Castro-Vincenzi et~al.(2024)Castro-Vincenzi, Khanna, Morales, and Pandalai-Nayar]{CastroVincenziEtAl2024}
Castro-Vincenzi, J., Khanna, G., Morales, N., and Pandalai-Nayar, N. (2024).
\newblock Weathering the storm: Supply chains and climate risk.
\newblock NBER Working Paper 32218, revised March 2025. doi:10.3386/w32218.

\bibitem[Ehime Prefecture(2026)]{EHIME2026}
Ehime Prefecture (2026).
\newblock Overview of economic and labor policy measures.
\newblock \url{https://www.pref.ehime.jp/page/8940.html}.

\bibitem[Eisenack(2014)]{Eisenack2014}
Eisenack, K. (2014).
\newblock The inefficiency of private adaptation to pollution in the presence of endogenous market structure.
\newblock \emph{Environmental and Resource Economics} 57(1), 81--99. doi:10.1007/s10640-013-9667-6.

\bibitem[Grames et~al.(2019)Grames, Grass, Kort, and Prskawetz]{GramesEtAl2019}
Grames, J., Grass, D., Kort, P.~M., and Prskawetz, A. (2019).
\newblock Optimal investment and location decisions of a firm in a flood risk area using impulse control theory.
\newblock \emph{Central European Journal of Operations Research} 27(4), 1051--1077. doi:10.1007/s10100-018-0532-0.

\bibitem[Ihara(2008)]{Ihara2008}
Ihara, R. (2008).
\newblock Transport costs, capital mobility and the provision of local public goods.
\newblock \emph{Regional Science and Urban Economics} 38(1), 70--80. doi:10.1016/j.regsciurbeco.2008.01.007.

\bibitem[Markusen et~al.(1993)Markusen, Morey, and Olewiler]{MarkusenMoreyOlewiler1993}
Markusen, J.~R., Morey, E.~R., and Olewiler, N.~D. (1993).
\newblock Environmental policy when market structure and plant locations are endogenous.
\newblock \emph{Journal of Environmental Economics and Management} 24(1), 69--86. doi:10.1006/jeem.1993.1005.

\bibitem[Markusen et~al.(1995)Markusen, Morey, and Olewiler]{MarkusenMoreyOlewiler1995}
Markusen, J.~R., Morey, E.~R., and Olewiler, N. (1995).
\newblock Competition in regional environmental policies when plant locations are endogenous.
\newblock \emph{Journal of Public Economics} 56(1), 55--77. doi:10.1016/0047-2727(94)01419-O.

\bibitem[Matsumoto(1998)]{Matsumoto1998}
Matsumoto, M. (1998).
\newblock A note on tax competition and public input provision.
\newblock \emph{Regional Science and Urban Economics} 28(4), 465--473. doi:10.1016/S0166-0462(98)00003-9.

\bibitem[Maurer and Walz(2000)]{MaurerWalz2000}
Maurer, B. and Walz, U. (2000).
\newblock Regional competition for mobile oligopolistic firms: Does public provision of local inputs lead to agglomeration?
\newblock \emph{Journal of Regional Science} 40(2), 353--375. doi:10.1111/0022-4146.00178.

\bibitem[Meng et~al.(2024)Meng, Ding, and Wang]{MengDingWang2024}
Meng, X., Ding, T., and Wang, H. (2024).
\newblock Urban resilience under local government competition: A new perspective on industrial resilience.
\newblock \emph{Cities} 155, 105409. doi:10.1016/j.cities.2024.105409.

\bibitem[Ministry of Economy, Trade and Industry(2026)]{METI2026}
Ministry of Economy, Trade and Industry, Japan (2026).
\newblock Industrial location policy.
\newblock \url{https://www.meti.go.jp/policy/local_economy/industrial_location/index.html}.

\bibitem[Ministry of the Environment(2026)]{MOE2026}
Ministry of the Environment, Japan (2026).
\newblock Manual for formulating regional climate change adaptation plans, revised March 31, 2026.
\newblock \url{https://www.env.go.jp/earth/earth/tekiou/page_00005.html}.

\bibitem[Otazawa et~al.(2016)Otazawa, Nakamura, Torio, and Koike]{OtazawaEtAl2016}
Otazawa, T., Nakamura, Y., Torio, K., and Koike, A. (2016).
\newblock Industrial location and vulnerability to natural hazards.
\newblock \emph{Journal of Japan Society of Civil Engineers, Ser. D3 (Infrastructure Planning and Management)} 72(1), 99--112. doi:10.2208/jscejipm.72.99.

\bibitem[Small and Medium Enterprise Agency(2026)]{SMEA_BCP}
Small and Medium Enterprise Agency, Japan (2026).
\newblock Business continuity guidelines: Identifying and selecting alternatives for business continuity.
\newblock \url{https://www.chusho.meti.go.jp/bcp/contents/level_b/bcpgl_03b_2_1.html}.

\bibitem[Walz and Wellisch(1996)]{WalzWellisch1996}
Walz, U. and Wellisch, D. (1996).
\newblock Strategic provision of local public inputs for oligopolistic firms in the presence of endogenous location choice.
\newblock \emph{International Tax and Public Finance} 3(2), 175--189. doi:10.1007/BF00399909.

\bibitem[Zhao et~al.(2026)Zhao, Yang, and Zhang]{ZhaoYangZhang2026}
Zhao, Y., Yang, Y., and Zhang, N. (2026).
\newblock Trading resilience for growth: Evidence from public flood protection in China.
\newblock SSRN working paper 6996377, posted June 25, 2026.
\end{thebibliography}

\end{document}
